\section{Related work}
\label{sec:related}

\paragraph{Classical calibration of head-mounted eye trackers.}
Linear and affine transforms~\citep{vadillo2015simple} provide a
closed-form baseline. Polynomial regression has been the workhorse for a
decade~\citep{cerrolaza2012polynomial,mardanbegi2018distribution}, with
quadratic forms repeatedly identified as the best
trade-off~\citep{ehinger2019eyelink,su2020locally}. Non-parametric
methods --- thin-plate spline, radial basis function with
multi-quadric kernels~\citep{scheel2016rbf}, and Gaussian process
regression~\citep{tripathi2017gpr} --- relax the global-parametric
assumption and capture local variation. In practice, a hybrid
\emph{similarity + RBF} (sim+RBF), in which a global affine fit is
followed by an RBF interpolation on the residual, is the
strongest single classical baseline on both JuDo1000 and our
self-collected data (see Section~\ref{sec:experiments}).

\paragraph{Neural and residual-based refinement.}
Several recent works augment classical pipelines with neural networks.
\citet{huang2025fgd} apply feature-guided de-drifting on EOG signals;
\citet{wang2017blind} learn a projection-recovery network for sensor
networks; \citet{thieu2025stochastic} use extended Kalman filtering for
gaze trajectory smoothing. The recent ``Multi Baseline Neural
Refinement'' approach we audit \citep{anon2026mbnr} feeds the
predictions of $K=7$ classical baselines into a $[64,32,16]$ MLP with
batch normalization and Gaussian noise injection, and reports
state-of-the-art accuracy on JuDo1000 and a self-collected dataset.
We show that the reported numbers depend on a single line of input-feature
construction that subtracts the regression target.

\paragraph{Online and implicit recalibration.}
Implicit calibration methods exploit natural viewing behaviour to keep
calibration current without explicit user
intervention~\citep{huang2019saccalib,muller2019mobile,wan2024selfcalib}.
Explicit short recalibration~\citep{lander2016time,hou2025xr} re-collects
a small number of fixations to refresh accuracy. Our online
shrinkage estimator falls into the explicit-but-very-short category,
collecting $K \le 18$ extra fixations.

\paragraph{Shrinkage estimation.}
Shrinking an empirical estimator toward a fixed target reduces variance
at the cost of a small bias and is well-studied in
statistics~\citep{james1961estimation,efron1977stein,ledoit2004wellconditioned}.
Our learned shrinkage MLP can be viewed as a non-parametric
generalization of a James-Stein estimator: instead of choosing a single
shrinkage scalar from theory, we predict a per-trial 2-D shrinkage from
sample statistics. The per-trial nature is central --- the bias varies
substantially across sessions, and the optimal $\lambda$ varies with
the noise level of the empirical mean, which depends on $K$.
